\documentclass[10pt]{article} 

\usepackage[a4paper,margin=2cm,left=2cm,right=2cm]{geometry} 
\usepackage[UKenglish]{babel}

\usepackage{amsmath}
\usepackage{amssymb}
\usepackage{amsthm}

\usepackage{multicol}

\newtheorem{theorem}{Theorem}
\newtheorem{lemma}{Lemma}
\newtheorem{corollary}{Corollary}[theorem]
\newtheorem{proposition}{Proposition}
\theoremstyle{definition}
\newtheorem{definition}{Definition}
\newtheorem{example}{Example}
\theoremstyle{remark}
\newtheorem{remark}{Remark}

\newcommand{\R}{\mathbb{R}}
\newcommand{\Z}{\mathbb{Z}}
\newcommand{\N}{\mathbb{N}}
\newcommand{\Q}{\mathbb{Q}}
% \DeclareMathOperator{\dom}{dom}

\title{Python Quick Reference}
\author{joel}
\date{}

\begin{document}

\maketitle

\section{SymPy}

Generally want to do \verb|from sympy import *| otherwise things get annoying fast. 
Also, using the REPL is actually quite helpful. 
Can do \verb|uv run --with sympy python| for a REPL environment with SymPy. 
\verb|init_printing()| will `pretty print' results in a way that is often more
readable than the regular form. 
 
\subsection{Basic operations}

Create symbols, expressions and evaluate constants:
\begin{verbatim}
>>> x = symbols('x')
>>> y = symbols('y')
>>> expr = (x + 1)**4       
>>> pi.evalf(10) 
3.141592654
\end{verbatim}
Note that \verb|expr| here is arbitrary, it's just a variable name. 

There are two ways to evaluate an expression at a particular value:
\begin{verbatim}
>>> expr.evalf(subs={x: 3})
256.000000000000
>>> expr.subs(x, 3).evalf()
256.000000000000
\end{verbatim}

For multiple substitutions, use a list of tuples, like:
\begin{verbatim}
>>> expr1 = (x + y*x + y**2)**3
>>> expr1.subs([(x, 1), (y, 2)]).evalf()
343.000000000000
\end{verbatim}

Factorise expressions: 
\begin{verbatim}
>>> expr2 = x**4 - 1
>>> expr2.factor()
(x - 1)*(x + 1)*(x**2 + 1)
\end{verbatim}

Simply expressions: 
\begin{verbatim}
>>> expr3 = (x**2 - 4) / (x - 2)
>>> expr3.simplify()
x + 2
>>> cancel(expr3)
x + 2
\end{verbatim}

\textbf{N.B Quotients:} using / is fine when symbols are involved, but when you have 
a quotient of two explicit numbers $a, b$, then use \verb|Rational(a, b)|. 

\subsection{Calculus} 

Find limits of expressions ($\infty$ is denoted by `oo'): 
\begin{verbatim}
    >>> limit(sin(x)/x, x, 0) 
    1
    >>> limit(1/x, x, oo)
    0
    >>> limit(1/x, x, 0, '+')
    oo
    >>> limit(1/x, x, 0, '-')
    -oo
\end{verbatim}


\end{document}
